\documentclass[11pt,a4paper]{article}
\usepackage[margin=23mm]{geometry}
\usepackage[T1]{fontenc}
\usepackage[utf8]{inputenc}
\usepackage{lmodern,microtype,amsmath,amssymb,booktabs,tabularx,array}
\usepackage{graphicx,xcolor,listings,enumitem}
\usepackage[section]{placeins}
\usepackage[colorlinks=true,linkcolor=blue!50!black,urlcolor=blue!50!black]{hyperref}
\usepackage{fancyhdr}
\pagestyle{fancy}
\fancyhf{}
\lhead{Statistical learning: worked notebook explanations}
\rfoot{\thepage}
\setlength{\headheight}{14pt}
\setlength{\parindent}{0pt}
\setlength{\parskip}{6pt}
\setlength{\emergencystretch}{2em}
\setlist{itemsep=3pt,topsep=4pt}
\lstset{language=Python,basicstyle=\small\ttfamily,breaklines=true,
 columns=fullflexible,keepspaces=true,showstringspaces=false,
 frame=single,rulecolor=\color{gray!35},backgroundcolor=\color{gray!4},
 commentstyle=\color{green!35!black},keywordstyle=\color{blue!65!black},
 aboveskip=8pt,belowskip=8pt}
\newcommand{\code}[1]{\texttt{\detokenize{#1}}}
\newcommand{\RSS}{\operatorname{RSS}}
\newcommand{\RMSE}{\operatorname{RMSE}}
\newcommand{\MAE}{\operatorname{MAE}}
\newcommand{\E}{\mathbb{E}}
\newcommand{\ind}[1]{\mathbf{1}\{#1\}}
\newcommand{\fig}[2]{\begin{figure}[htbp]\centering
\includegraphics[width=\linewidth,height=0.39\textheight,keepaspectratio]{figures/#1}
\caption{#2}\end{figure}}
\newcommand{\notebooktitle}[3]{\begin{center}{\LARGE\bfseries #1\par}\vspace{5pt}
{\large Problem statement and worked resolution\par}\vspace{5pt}
Companion to \code{#2}\par\end{center}
\textbf{How to use this explanation.} #3 Section numbers follow the notebook's numbered sections where present. Numerical values are rounded and refer to the recorded run, not universal performance claims.}
\newcommand{\methods}[1]{\subsection*{Main Python calls}#1}

\begin{document}
\notebooktitle{3C: OJ purchases and score-based decisions}{03_C_oj_score_thresholds.ipynb}{This resolution explains the prediction setting, score-threshold search, and why a policy can lower cost while making more classification mistakes.}
\section*{Problem statement and scope}
The textbook OJ data contain 1,070 purchases: 653 CH and 417 MM. Encode MM as class 1 and CH as 0.
Predict from PriceCH, PriceMM, DiscCH, DiscMM, SpecialCH, and SpecialMM. The first four describe prices and discounts,
and the last two promotion indicators. The target Purchase never appears among the inputs.
This is a specified restricted-feature benchmark, not a claim that the six variables contain every useful predictor.

The SVM learns a decision score $f(x)$. A policy predicts MM when $f(x)\ge t$.
The costs are assumed: predicting MM for a CH purchase costs 1; predicting CH for an MM purchase costs 3.
They could represent asymmetric consequences in a teaching scenario, but are not estimated from the records.
We first select an SVM, then select its threshold using a different data partition.

\section{Load data and reserve partitions}
The notebook uses the public textbook repository URL with a 20-second network timeout.
This avoids requiring local course files in Colab. Shape, label-count, and missing-value checks verify the expected data.
The \href{https://intro-stat-learning.github.io/ISLP/datasets/OJ.html}{OJ data documentation} provides variable context.

Stratification with seed 23 produces 642 training, 214 validation, and 214 test rows.
Training contains 392 CH and 250 MM records. Price columns have substantially different numerical ranges
from discount and binary promotion columns; each pipeline standardizes the six columns using fitting rows only.
The same stores occur in every partition and their week ranges overlap.
Consequently this random-split result concerns similarly sampled records. It does not show how a model would
perform in an unseen store or in future weeks, and row-index separation alone cannot establish independence of repeated purchases.

\section{Select model settings without test information}
For logistic, linear SVM, and RBF SVM, tune C in $\{0.01,0.1,1,10\}$;
RBF also tunes gamma in $\{0.01,0.1,1,10\}$. Each pipeline includes its own fitting scaler.
Five training folds have slightly different sizes because 642 is not divisible by five:
held-out folds have 128 or 129 rows, with the remaining 514 or 513 used to fit.

Mean balanced accuracy gives equal weight to CH recall and MM recall. It does not minimize the later cost directly.
That separation is deliberate: fitting a useful score and choosing a cost-sensitive cutoff answer different questions.
\begin{center}\begin{tabular}{lrrr}\toprule
Family&Selected C&Gamma&CV balanced accuracy\\\midrule
Logistic&0.1&---&0.622797\\
Linear SVM&1&---&0.582513\\
RBF SVM&1&0.1&0.630041\\\bottomrule\end{tabular}\end{center}
The selected SVM is the RBF model. Logistic regression remains a tuned comparison, not a compulsory loser.
All three models are refitted on the full training partition and left unchanged thereafter.
The chosen model's class order is checked as $[0,1]$, so positive SVM scores favor MM.

\section{Inspect scores and search every decision pattern}
There are only 64 distinct validation scores, giving 65 possible threshold decision patterns.
Repeated or equivalent input records can receive identical scores. A threshold cannot split a tied-score group
into different labels without introducing another decision rule.

Let sorted distinct scores be $s_1<\cdots<s_k$. Under the rule $f\ge t$, use cutoffs
\[-\infty,s_2,s_3,\ldots,s_k,+\infty.\]
The first predicts MM for every record, the last CH for every record, and each intermediate cutoff excludes
one more lowest-score group. This covers all patterns without arbitrary fine grid spacing.
The exact magnitude of a representative cutoff is less important than the group of records it classifies positive.

For every pattern, count false positives FP and false negatives FN and minimize $FP+3FN$.
Division by 214 converts this to mean cost but does not change its minimum. Among equal-cost cutoffs, use the largest,
which favors fewer positive predictions. Including extreme policies prevents overlooking an optimum that always predicts one class.

The selected cutoff is approximately $-1.001416$, with validation cost $117/214\approx0.546729$.
Always predicting MM has cost 130; always predicting CH has cost $3\times84=252$ on validation.
The plot displays finite cutoffs only, but the extremes remain in the optimization and printed table.
\fig{03_C_oj_score_thresholds_1.png}{Validation score distributions and threshold costs. The zero cutoff is a default score convention; the selected cutoff reflects this sample and the specified costs.}
The cost curve is piecewise constant between score values. Connecting or stepping between plotted values is a
visualization convention, not evidence that costs vary smoothly with every tiny threshold movement.

\section{Why the probability threshold 1/4 does not apply to scores}
For a true conditional MM probability $p$, the expected cost of predicting MM is $1\times(1-p)$,
and the cost of predicting CH is $3\times p$. Prefer MM when $1-p\le3p$, hence $p\ge1/4$.
This derivation concerns a probability scale. A raw SVM score can be negative and has no such scale.
The selected score near $-1$ is therefore not inconsistent with the probability argument.

Applying $\sigma(f)=1/(1+e^{-f})$ without learning a calibration mapping would merely transform numbers;
it would not establish that they equal purchase risks. A strictly increasing transform preserves score ordering
and the possible threshold policies, but changes the numerical cutoff. For example, transforming both scores
and their selected cutoff monotonically can preserve every decision without improving calibration at all.

After threshold selection, retain the fitted model and scaler. Refitting them on training plus validation
could change scores and invalidate the cutoff, even if the new model had similar overall accuracy.
The focused experiment does not perform calibration or report probability losses for these raw scores.

\section{Final test results and the cost--accuracy tradeoff}
At the default zero-score policies:
\begin{center}\begin{tabular}{lrrr}\toprule
Model&Accuracy&Balanced accuracy&ROC AUC\\\midrule
Training-majority&0.6122&0.5000&0.5000\\
Logistic&0.6402&0.5935&0.6509\\
Linear SVM&0.6495&0.5769&0.6480\\
RBF SVM&0.6776&0.6197&0.6731\\\bottomrule\end{tabular}\end{center}
ROC AUC evaluates ranking from real scores, not the selected cutoff, and does not require calibrated probabilities.
It is distinct from both balanced accuracy and operational cost.
For the selected RBF SVM:
\begin{center}\begin{tabular}{lrrrrrr}\toprule
Policy&TN&FP&FN&TP&Accuracy&Cost/record\\\midrule
Zero score&115&16&53&30&0.6776&0.8178\\
Validation cutoff&18&113&1&82&0.4673&0.5421\\\bottomrule\end{tabular}\end{center}
The changed cutoff prevents 52 misses while adding 97 false positives. The cost change is
\[97-3\times52=-59,\]
so mean cost falls by $59/214\approx0.2757$ despite the increase in the ordinary error count.
The selected policy has precision $82/195\approx0.4205$ and recall $82/83\approx0.9880$.
This is a high-recall, many-false-alarm policy, not a universally better classifier.

For scale, always predicting MM on this test would cost $131/214\approx0.6121$;
always predicting CH would cost $249/214\approx1.1636$. These transparent reference calculations show what
an asymmetric-cost improvement means beyond a comparison with the default zero cutoff.
They do not trigger reselection after test evaluation. The validation-selected policy remains fixed regardless
of whether any reference happens to look better in another test sample.

\section*{Documentation and practical advice}
The \href{https://scikit-learn.org/1.8/modules/svm.html}{SVM documentation} distinguishes decision scores from
probability estimation. Check \code{classes_} when translating score signs, and keep the complete preprocessing
pipeline when saving a policy. A threshold value without its fitted score function is not a reproducible classifier.

\begin{itemize}
\item Name the positive event before calculating FP, FN, precision, recall, or cost; reversing labels changes the interpretation.
\item Search tied scores as groups and include the two constant policies. An arbitrary numeric grid may omit relevant decisions.
\item A cost minimum fitted on validation has selection noise. Evaluate it once on the reserved test.
\item Never interpret a price coefficient, support vector, or selected cutoff as evidence of a causal price intervention.
\item For a future-store objective, change the evaluation design first; better random-split accuracy cannot answer that question on its own.
\end{itemize}
\end{document}
